% =========================================================================
%  Abstract and keywords (IEEE: no math, symbols or citations)
% =========================================================================
\begin{abstract}
The sawtooth wave of synthesizers contains more high
frequencies than a digital signal can hold. When software computes it
sample by sample, the excess folds back as unwanted tones
called aliasing. This paper reviews three families
of antialiasing oscillators and compares a naive sawtooth, the polynomial
bandlimited step (PolyBLEP) and the second-order differentiated parabolic
waveform (DPW-2) with an ideal reference. PolyBLEP lowers the
aliasing by a median of 18 dB and DPW-2 by 11 dB, at about 1.3 to 1.5
times the cost of the naive oscillator. Two musical examples show that
aliasing matters most in high melodies.
\end{abstract}

\begin{IEEEkeywords}
aliasing, sound synthesis, oscillators, PolyBLEP
\end{IEEEkeywords}
